\label{sec:results}

\subsection{Experimental Design}

Quantitative results for M.4 commuting shed weights are deferred to Phase~2,
pending download of LODES OD data for the target states.\footnote{Phase~2
will be triggered after LODES OD files are downloaded.
The data URL pattern is
\texttt{https://lehd.ces.census.gov/data/lodes/LODES8/\{state\}/od/\{state\}\_od\_main\_JT00\_2020.csv.gz}.
Files are downloaded via \texttt{bisect fetch --type lodes-od --year 2020}.}
This section describes the experimental protocol governing Phase~2 execution.

\paragraph{States.}
Two states are selected to test distinct aspects of the commuting shed signal:

\begin{itemize}
  \item \textbf{Wisconsin ($k=8$).}
    WI provides the suburban clustering test.
    The Milwaukee metro is characterized by distinct commuting sheds:
    the western suburbs (Waukesha County) commute primarily to Milwaukee
    and the western suburban employment corridor (Brookfield, Wauwatosa,
    Menomonee Falls, Pewaukee); the northern suburbs (Ozaukee, Washington
    counties) commute primarily to Milwaukee and the northern suburban
    corridor (Glendale, Brown Deer, Menomonee Falls).
    Both sheds overlap at downtown Milwaukee as a common destination.
    The prediction: tracts in these suburban sheds co-assign to the same
    district at ${\geq}70\%$ under commuting-shed weights, versus
    ${\leq}50\%$ under geographic weights.
  \item \textbf{North Carolina ($k=14$).}
    NC provides the polycentric separation test.
    The Research Triangle (Wake + Durham + Orange + Chatham counties) and
    the Charlotte metro (Mecklenburg + Union counties) are two distinct
    functional economic areas separated by 160 miles of Piedmont geography.
    Under geographic weights, a plan that bundles some Piedmont tracts with
    the Triangle or with Charlotte can produce cross-metro district assignments.
    Under commuting-shed weights, the disjoint destination distributions of
    Triangle-commuting and Charlotte-commuting tracts should prevent such
    cross-metro assignments.
\end{itemize}

All experiments use:
\begin{itemize}
  \item Algorithm: \texttt{standard-bisect} structure
  \item Seed: 42 (single seed; multi-seed variance designated Phase~3)
  \item Blend factor: $\alpha = 0.5$
  \item Year: 2020 decennial census population and 2020 LODES OD
  \item Comparison baseline: \texttt{--weights-override geographic}
    on identical plan configuration
\end{itemize}

\subsection{Primary Metric: Suburban Clustering (Wisconsin)}

The WI primary metric is the \emph{Milwaukee suburban co-assignment rate}:
the fraction of tract pairs $(h_1, h_2)$ drawn from the Milwaukee
suburban counties (Waukesha, Ozaukee, Washington) that are assigned to
the same district.

\begin{equation}
  \mathrm{SCR} = \frac{\#\{(h_1, h_2) : h_1, h_2 \in \mathcal{S},\, \pi(h_1) = \pi(h_2)\}}
  {\#\{(h_1, h_2) : h_1, h_2 \in \mathcal{S}\}}
  \label{eq:scr}
\end{equation}
where $\mathcal{S}$ is the set of Milwaukee suburban tracts and
$\pi(h)$ is the district assignment of tract $h$.

The Phase~2 prediction:
$\mathrm{SCR}_{\mathrm{commute}} \geq 0.70$ and
$\mathrm{SCR}_{\mathrm{geographic}} \leq 0.50$.\textsuperscript{$\dagger$}

\subsection{Primary Metric: Polycentric Separation (North Carolina)}

The NC primary metric is the \emph{cross-metro co-assignment rate}:
the fraction of tract pairs $(h_1, h_2)$ where $h_1$ is in the
Research Triangle metro and $h_2$ is in the Charlotte metro that are
assigned to the same district.

\begin{equation}
  \mathrm{CMR} = \frac{\#\{(h_1, h_2) : h_1 \in \mathcal{T},\, h_2 \in \mathcal{C},\, \pi(h_1) = \pi(h_2)\}}
  {\#\{(h_1, h_2) : h_1 \in \mathcal{T},\, h_2 \in \mathcal{C}\}}
  \label{eq:cmr}
\end{equation}
where $\mathcal{T}$ is the set of Research Triangle metro tracts and
$\mathcal{C}$ is the set of Charlotte metro tracts.

The Phase~2 prediction:
$\mathrm{CMR}_{\mathrm{commute}} < 0.05$ and
$\mathrm{CMR}_{\mathrm{geographic}} \leq 0.15$.\textsuperscript{$\dagger$}

A cross-metro co-assignment rate below 5\% under commuting-shed weights
would confirm that the weighted Jaccard similarity correctly identifies
the Triangle and Charlotte as distinct functional economic areas that should
not be bundled into the same district.

\subsection{Performance Validation}

A third metric validates the computational cost claim:

\begin{quote}
  \textit{Wall time for weighted Jaccard computation across all NC tract
  adjacency edges: $<60$ seconds on 8 workers.}
\end{quote}

This is measured by the \textsc{bisect} internal timer for the weight
construction phase, reported in the \texttt{bisect state} diagnostic
output.
The prediction follows from the $O(n \times 2\bar{W})$ complexity analysis
in Section~\ref{sec:algorithm}: with $n \approx 7{,}000$ NC adjacency edges
and $\bar{W} \approx 20$, approximately 280,000 operations are required,
a trivially fast computation even on a single thread.\textsuperscript{$\dagger$}

\subsection{Phase 2 Timeline}

All quantitative results in this section are marked ${}^{\dagger}$Phase~2
and will be populated when the following prerequisites are met:

\begin{enumerate}
  \item LODES OD files for NC and WI downloaded via
    \texttt{bisect fetch --type lodes-od --year 2020}.
  \item Block-to-tract aggregation completed and sparse commuting
    distribution matrices stored in Parquet format.
  \item L0 test invariants verified for the aggregated distributions
    (Section~\ref{sec:algorithm}).
  \item Single-seed redistricting runs completed for both states under
    both \texttt{commuting-shed} and \texttt{geographic} weight modes.
  \item SCR, CMR, and wall time metrics computed from plan output.
\end{enumerate}
